\subsection{Алгоритмические задачи}

\begin{taskbox}[Плавная пульсация яркости красного канала (PWM-модуляция)]
Необходимо реализовать алгоритм плавного увеличения и уменьшения яркости красного канала от значения \verb|0.0| до \verb|1.0| и обратно. Изменение яркости должно происходить малыми приращениями с использованием таймера.

\textbf{Ожидаемый результат:} вся светодиодная линейка плавно пульсирует красным цветом.
\end{taskbox}

\begin{taskbox}[Стробоскопический режим (белые вспышки)]
Требуется реализовать быстрое периодическое включение и выключение всей светодиодной линейки белым цветом \verb|{1,1,1}| с малым временным интервалом.

\textbf{Ожидаемый результат:} частые и чёткие белые вспышки, создающие эффект стробоскопа.
\end{taskbox}

\begin{taskbox}[Отображение чётности индекса (шахматный узор)]
Необходимо реализовать статический шаблон индикации, при котором светодиоды с чётными индексами светятся красным цветом \verb|{1,0,0}|, а с нечётными — синим \verb|{0,0,1}|.

\textbf{Ожидаемый результат:} устойчивый чередующийся цветовой узор по всей линейке.
\end{taskbox}

\begin{taskbox}[Инвертируемая полицейская мигалка]
Требуется реализовать алгоритм, при котором цвета чётных и нечётных светодиодов периодически меняются местами с заданным интервалом времени.

\textbf{Ожидаемый результат:} регулярное чередование цветовых паттернов, имитирующее работу полицейской световой сигнализации.
\end{taskbox}

\begin{taskbox}[Генерация цветового шума]
Необходимо реализовать алгоритм, который с заданной периодичностью устанавливает для каждого светодиода случайные значения компонент RGB в диапазоне от 0 до 1.

\textbf{Ожидаемый результат:} динамично изменяющаяся цветовая картина без видимой закономерности.
\end{taskbox}
